% Literature-grounded draft, revised following the author's clarification.
% Study material: 304 stainless steel, represented by Fe--Cr--Ni--Si.
% All five cited papers have local full texts; see literature/README.md.

Grain-boundary oxidation in 304 stainless steel (304SS) exposed to
pressurized water reactor (PWR) primary water depends on the chemical
changes introduced by prior irradiation. Deng and co-workers found that
increasing the proton irradiation dose produced deeper intergranular
oxidation in 304 nuclear-grade stainless steel after water exposure,
accompanied by a decrease in the Cr content of the intergranular oxide
\parencite{deng2017}. These observations connect irradiation history to
both oxide chemistry and penetration depth. Establishing a quantitative
relationship between them requires a description of how the changing
oxide composition affects oxygen transport and, consequently, the
advance of the oxidation front.

Radiation-induced segregation (RIS) establishes the local alloy
composition from which the intergranular oxide develops. Coupling between
solute and irradiation-generated point-defect fluxes redistributes
alloying elements near grain boundaries. Measurements on irradiated
304 nuclear-grade stainless steel reveal Cr depletion together with Ni
and Si enrichment \parencite{deng2017,deng2021apt}. Subsequent oxidation
selectively consumes alloying elements, changing both the oxide composition
and the adjacent metal, where Ni enrichment has also been observed
\parencite{deng2021apt}. The lower Cr content of the oxide formed on
irradiated boundaries has been associated with reduced resistance to
oxygen diffusion \parencite{deng2017}. Oxide protectiveness therefore
has a direct transport meaning: an oxide that restricts oxygen diffusion
more strongly limits the oxygen available for further oxidation.

Si oxidation and dissolution also alter this transport resistance by
changing the structure through which oxygen moves. Using model alloys
with systematically varied Cr and Si contents, Du and co-workers linked
increased Si content to oxide degradation and greater cracking
susceptibility in hydrogenated water. They attributed this behavior to
dissolution of Si-containing oxide and the resulting porous structure,
which facilitates oxygen transport \parencite{du2026}. Related microscopy
on proton-irradiated 316L identified Si enrichment at oxide tips and
porosity behind them, consistent with preferential Si oxidation followed
by dissolution \parencite{kuang2022}. These studies provide a mechanistic
basis for including dissolution-generated free volume in the transport
description. Together with the dependence on Cr content, they motivate
representing the evolving oxide composition and free volume through an
effective oxygen diffusivity. This connects the segregation of multiple
elements to a common quantity governing oxygen transport.

Models coupling solute and point-defect transport can predict the
composition redistribution produced by irradiation. For example,
Kadambi and co-workers developed a model for segregation at grain
boundaries and dislocation structures in austenitic Fe--Cr--Ni
\parencite{kadambi2025}. Connecting such predictions to oxidation
kinetics requires tracking how the segregated alloy supplies the
interfacial reactions, how those reactions change the oxide, and how
the resulting diffusivity affects subsequent oxygen delivery. The
quantity to be predicted is thus the oxidation response arising from
the evolving chemical state. Segregation profiles, oxide composition,
interfacial Ni enrichment, and penetration depth provide complementary
observables for assessing this coupling.

Here, we develop a coupled RIS and grain-boundary oxidation model for
304SS, represented by the Fe--Cr--Ni--Si system. The model links
defect-mediated solute redistribution and selective interfacial
consumption to an effective oxygen diffusivity, \(D_{\mathrm{eff}}\),
determined by the evolving oxide composition and free volume. Within
this description, increasing the Cr content of the oxide lowers the
effective diffusivity, whereas Si oxidation followed by complete
dissolution contributes free volume that enhances transport. The same
\(D_{\mathrm{eff}}\) enters both oxygen transport along the grain boundary
and oxygen delivery across the oxide to the metal interface. The model
therefore connects RIS to oxidation kinetics through the resulting change
in oxygen transport, with metal consumption feeding back on the local
composition. The complete-dissolution treatment is a simplification
motivated by the mechanism discussed by Du and co-workers
\parencite{du2026}. We examine the dose and exposure-time dependence of
oxide penetration under sequential irradiation and water exposure, using
separate calibration of the RIS and oxidation components and a common
set of oxidation parameters across doses. The framework also permits
predictions of concurrent RIS and oxidation within the same transport
description.
% The exposure-time dataset comes from a separate paper, as clarified by
% the author. Add its citation once identified; do not attribute it to
% Deng 2017 or assume it is Wang 2023.
